Low resolution can damage staff lines, symbols, digits, and text in a digitized music score.
Super-resolution can generate a larger image, but models trained on natural images may introduce
visually plausible detail that was not present in the document. This Bachelor's Thesis studies
when these techniques add value and which safeguards are required for professional use in
libraries, archives, preservation, consultation, or editorial preparation.

Sheet Music Benchmark (SMB) was audited and an untouched sample was reserved for the primary
comparison. It used 64 pages from 64 documentary groups and six synthetic degradations: two
and four times upscaling crossed with clean, moderate, and strong severity. Blur was normalized by
staff-line spacing so that severity remained comparable across pages. Bicubic interpolation, the
EDSR residual network, and lightweight SwinIR were compared using luminance PSNR and SSIM,
paired intervals by group, inference time, and a preselected qualitative review of notation
failures. The reproducible implementation was developed in Python with PyTorch, OpenCV, and
\texttt{uv}; GPU execution used one Tesla T4 device on Kaggle.

After preserving the primary comparison, a secondary EDSR fine-tuning study was performed.
Project-defined roles comprised 45 training, 13 validation, and 20 test groups, with no identifier
overlap or reuse of the 64 primary groups. Groups can be movements of the same composition:
musical independence is not guaranteed, and intervals are descriptive and conditional on this
grouping. Adapted EDSR exceeded the official model in luminance PSNR and SSIM under all six
conditions, with group-paired intervals that did not cross zero. This supports adaptation on
held-out SMB documents under matched degradations, without certifying generalization to unseen
complete compositions or musical correctness.

As a professional extension, the adapted model was applied without retuning to twelve works from
an external music archive under the same degradations. Intervals favored the adapted model in all
six PSNR comparisons and five of six SSIM comparisons. Eight derivatives were accepted for
consultation, three with reservations, and one was rejected; a reversible local demonstrator was
also implemented. This supports cross-collection transfer under synthetic LR, not restoration of
naturally acquired LR inputs or a universal safety rate.

EDSR and SwinIR outperformed bicubic interpolation on both primary metrics in all six conditions.
However, SwinIR clearly exceeded EDSR only for x2 clean, x2 moderate, and x4 clean inputs, while
requiring 5.2 to 8.5 times more inference time. Visual review found altered symbols, corrupted text,
and curved or textured strokes even when average fidelity improved. EDSR therefore provides the
most defensible fidelity--cost trade-off for assisted access copies, whereas SwinIR should be
reserved for conditions in which its additional gain justifies the cost. None of the evaluated
methods is validated as semantic restoration, an archival replacement, or an autonomous input to
optical music recognition. The original must be retained and every consequential derivative must
be checked by a human reviewer.
